%% ma: L12-B15-0025
%% lop: 12
%% chuong: 5
%% bai: 15
%% dang: 12.15.20
%% loai: TN4
%% muc: VD
%% dap_an: "D"
%% nguon: "(NBV)-10 ĐỀ KIỂM TRA HK2-TOÁN 12-DUNG CHUNG 3 SÁCH.pdf (H04, Đề số 4), Phần 1 Câu 3"
%% kien_thuc: "Bài 8 (tích có hướng), Bài 14 (vectơ pháp tuyến), Bài 15 (đường thẳng cắt và vuông góc với đường thẳng cho trước, nằm trong mặt phẳng)"
%% trang_thai: cho-duyet
%% ly_do: "Dạng mới đề xuất 12.15.20 (Đường thẳng nằm trong mặt phẳng, cắt và vuông góc với đường thẳng cho trước (tích có hướng hai vectơ)), chờ giáo viên duyệt."
%% ghi_chu: "Đáp án gốc D đúng."
\begin{ex}
Trong không gian $Oxyz$, cho đường thẳng $\Delta:\dfrac{x}{1}=\dfrac{y+1}{2}=\dfrac{z-1}{1}$ và mặt phẳng $(P):x-2y-z+3=0$. Đường thẳng $d$ nằm trong $(P)$ đồng thời cắt và vuông góc với $\Delta$ thì $d$ có phương trình là
\choice{$\begin{cases}x=1+2t\\y=1-t\\z=2\end{cases}$}{$\begin{cases}x=-3\\y=-t\\z=2t\end{cases}$}{$\begin{cases}x=1+t\\y=1-2t\\z=2+3t\end{cases}$}{\True $\begin{cases}x=1\\y=1-t\\z=2+2t\end{cases}$}
\loigiai{$\Delta:x=t,\ y=-1+2t,\ z=1+t$. Giao điểm $M=\Delta\cap(P)$: $t-2(-1+2t)-(1+t)+3=0\Leftrightarrow 4-4t=0\Leftrightarrow t=1$, nên $M(1;1;2)$.

$\overrightarrow{n_P}=(1;-2;-1)$, $\overrightarrow{u_\Delta}=(1;2;1)$. Vì $d\subset(P)$ và $d\perp\Delta$ nên $\overrightarrow{u_d}=\left[\overrightarrow{n_P},\overrightarrow{u_\Delta}\right]=(0;-2;4)$, chọn $(0;-1;2)$.

Phương trình $d$: $x=1$, $y=1-t$, $z=2+2t$. Chọn D.}
\end{ex}
